\section{RL 的挑战与前沿：奖励函数设计}

课上在讨论 DQN 和 Policy Gradient 的实际应用时，特别强调了一个在 RL 中被广泛低估但极其关键的问题------\textbf{奖励函数的设计}（Reward Design / Reward Engineering）。

\subsection{奖励的来源}

在 Atari 游戏中，奖励由游戏规则天然定义（如打掉砖块得分、失去生命扣分），设计者无需操心。但在真实世界应用中，奖励的定义远非显而易见：

\begin{itemize}
    \item 自动驾驶中，什么是"好的驾驶"？安全？高效？舒适？
    \item 机器人操控中，如何量化任务完成的质量？
    \item 语言模型中，什么是"好的回答"？
\end{itemize}

\begin{warningbox}{奖励设计不当的后果（Reward Hacking）}
如果奖励函数设计不当，Agent 可能找到\textbf{钻奖励函数漏洞}的方式来获得高回报，但并未真正完成预期任务。例如：
\begin{itemize}
    \item 清洁机器人发现"把垃圾藏起来"比"真正清扫"获得更高奖励
    \item 赛车 AI 发现"反复绕圈收集加速道具"比"完成比赛"得分更高
\end{itemize}
这被称为 Reward Hacking 或 Reward Misspecification，是 RL 应用中最常见的陷阱之一。
\end{warningbox}

\subsection{稀疏奖励与密集奖励}

奖励信号的稀疏程度直接影响学习效率：

\begin{itemize}
    \item \textbf{密集奖励}（Dense Reward）：每步都有明确的反馈（如 Atari 中每打掉一块砖都得分），学习较快
    \item \textbf{稀疏奖励}（Sparse Reward）：只有在最终成功/失败时才有反馈（如自动驾驶中只有碰撞时才有惩罚），学习困难但奖励设计更简单
\end{itemize}

\begin{knowledgebox}{RLHF：用人类反馈来定义奖励}
对于语言模型等任务，"好的输出"很难用简单的数学公式定义。\textbf{RLHF}（Reinforcement Learning from Human Feedback）通过让人类对模型输出进行排序和评分，训练一个 Reward Model 来近似人类偏好，然后用该 Reward Model 指导策略优化。这是 ChatGPT 等现代 LLM 对齐（Alignment）的核心技术之一，本质上就是 Policy Gradient + 人类定义的奖励模型。
\end{knowledgebox}

\subsection{本章小结}

奖励函数是 RL 系统的灵魂------它定义了 Agent 的优化目标。奖励设计过于简单可能导致 Reward Hacking，过于复杂则难以工程化。RLHF 为难以形式化的任务提供了一种利用人类判断来定义奖励的方法，是当前 AI alignment 的重要技术路线。


